\chapter{Теоретические сведения}
\label{ch:theory}

\section{Облачное хранилище и Seafile}

\textbf{Облачное хранилище} (cloud storage) --- модель хранения данных, при которой
файлы размещаются на удалённых серверах и доступны через сеть. Ключевые преимущества:
централизованное управление, синхронизация между устройствами, совместный доступ.

\textbf{Seafile} --- открытое ПО с открытым исходным кодом (лицензия AGPLv3 / Apache~2.0).
Основные компоненты:
\begin{itemize}
  \item \textbf{Seafile Server} --- демон хранения и синхронизации файлов, работает
        на порту~\texttt{8082};
  \item \textbf{Seahub} --- веб-приложение на Django, GUI и REST~API, порт~\texttt{8000};
  \item \textbf{Seafile Client} --- настольный клиент (Windows / macOS / Linux) и
        мобильные приложения для синхронизации.
\end{itemize}

\section{Стек технологий}

\textbf{MariaDB} --- реляционная СУБД, совместимая с MySQL, используется как бэкенд
для трёх баз данных Seafile: \texttt{ccnet-db}, \texttt{seafile-db}, \texttt{seahub-db}.

\textbf{Nginx} --- высокопроизводительный веб-сервер и обратный прокси.
В данной работе Nginx принимает HTTP-запросы на порту~\texttt{80}
и проксирует их на \texttt{127.0.0.1:8000} (Seahub):
\begin{lstlisting}
location / {
    proxy_pass http://127.0.0.1:8000;
    proxy_set_header Host $host;
    proxy_set_header X-Real-IP $remote_addr;
}
\end{lstlisting}

\textbf{Python~3 / pip} --- платформа для выполнения Seahub. Необходимые
Python-пакеты устанавливаются командой \texttt{pip3 install}:
\texttt{django}, \texttt{Pillow}, \texttt{mysqlclient}, \texttt{sqlalchemy},
\texttt{pycryptodome} и др.

\section{Схема взаимодействия компонентов}

После полной настройки клиентская машина Desktop обращается к серверу по имени
\texttt{seafile} (разрешается через BIND9 на gateway). Nginx принимает соединение
и перенаправляет запрос в Seahub. Seahub взаимодействует с Seafile-сервером
через CCNET~API; данные хранятся в директории \texttt{/opt/seafile/seafile-data}.
Информация о библиотеках, пользователях и правах хранится в MariaDB.

\section{DNS-запись для seafile}

Сервер \texttt{seafile} получает IP~\texttt{192.168.\cfgLabStudentN.4} и
регистрируется в прямой зоне BIND9 на \texttt{gateway}:
\begin{lstlisting}
seafile  IN  A  192.168.N.4
\end{lstlisting}
Клиент Desktop после этого может обращаться к серверу по короткому имени
\texttt{seafile} или по FQDN \texttt{seafile.STUDENT.GROUP.local}.

\endinput
